% ============================================================
\section{Implementation}
\label{sec:implementation}
% ============================================================

\kmerdet{} is written entirely in Rust (edition 2021) and compiles to a single
statically-linked binary with no runtime dependencies beyond the operating system
and (optionally) a system libjellyfish.  This section describes the key design
decisions, the module architecture, and how Rust's type system encodes domain
knowledge throughout the codebase.

\subsection{Why Rust}
\label{sec:why-rust}

The choice of Rust~\cite{matsakis2014rust,rustlang} over C++, Python, or Go
was driven by four properties that are particularly valuable for clinical
bioinformatics software:

\textbf{Memory safety without garbage collection}.  The borrow checker guarantees
at compile time that there are no use-after-free, double-free, or data-race
bugs.  In a tool that processes patient genomic data, silent memory corruption
errors are unacceptable.  The borrow checker eliminates this class of defect
entirely without the unpredictable pause times of a garbage collector, which
matters for deterministic performance in pipeline contexts.

\textbf{Zero-cost abstractions}.  Traits, generics, and iterators in Rust
compile to code as efficient as hand-written C.  The \texttt{KmerDatabase} trait
(Section~\ref{sec:trait-abstractions}) has zero runtime overhead: the compiler
monomorphises each concrete implementation, producing specialised code paths
without virtual dispatch.

\textbf{Native parallelism}.  The Rayon data-parallelism library~\cite{rayon}
provides safe parallel iteration through ownership semantics: calling
\texttt{par\_iter()} on a target vector processes all targets in parallel with
work-stealing load balancing, and the borrow checker enforces that no target
accesses shared mutable state (the jellyfish database is read-only and therefore
safe to share across threads without locking).

\textbf{Single-binary deployment}.  Clinical bioinformatics environments often
prohibit internet-connected package managers or complex runtime installations.
A single statically-linked binary with no Python interpreter, no virtual
environments, and no dynamic library dependencies can be deployed by copying a
single file, dramatically reducing deployment friction.

The primary trade-off is a steeper learning curve and longer initial development
time compared to Python.  For a tool intended for long-term use in clinical
pipelines, this investment is justified by the safety and performance guarantees.

\subsection{Module Architecture}
\label{sec:architecture}

\begin{figure}[t]
  \centering
  \begin{tikzpicture}[
    every node/.style={font=\small},
    mod/.style={rectangle, draw=gray!60, fill=gray!8, minimum height=0.8cm,
                minimum width=2.4cm, align=center},
    cli/.style={rectangle, draw=blue!50, fill=blue!8, minimum height=0.8cm,
                minimum width=2.4cm, align=center},
    arrow/.style={-{Latex[length=2mm]}, gray!60}
  ]

  % Top row: CLI layer
  \node[cli] (main) at (0, 5) {\texttt{main.rs}};
  \node[cli] (cli) at (3, 5) {\texttt{cli/}};
  \node[cli] (config) at (6, 5) {\texttt{config.rs}};

  % Middle row: Algorithm layer
  \node[mod] (kmer) at (-4.5, 3) {\texttt{kmer/}};
  \node[mod] (jf) at (-1.5, 3) {\texttt{jellyfish/}};
  \node[mod] (seq) at (1.5, 3) {\texttt{sequence/}};
  \node[mod] (walker) at (4.5, 3) {\texttt{walker/}};

  % Lower row
  \node[mod] (graph) at (-3, 1) {\texttt{graph/}};
  \node[mod] (variant) at (0, 1) {\texttt{variant/}};
  \node[mod] (filter) at (3, 1) {\texttt{filter/}};
  \node[mod] (output) at (6, 1) {\texttt{output/}};

  % Connections
  \draw[arrow] (main) -- (cli);
  \draw[arrow] (cli) -- (config);
  \draw[arrow] (cli) -- (walker);
  \draw[arrow] (jf) -- (walker);
  \draw[arrow] (kmer) -- (jf);
  \draw[arrow] (seq) -- (walker);
  \draw[arrow] (walker) -- (graph);
  \draw[arrow] (graph) -- (variant);
  \draw[arrow] (variant) -- (filter);
  \draw[arrow] (variant) -- (output);

  \end{tikzpicture}
  \caption{Module dependency graph for \kmerdet{}.  Blue nodes are the CLI layer;
    grey nodes are the algorithmic library.  Arrows indicate ``uses'' relationships.}
  \label{fig:module-graph}
\end{figure}

The module structure (Figure~\ref{fig:module-graph}) separates concerns cleanly:

\begin{description}
  \item[\texttt{kmer/}] 2-bit k-mer encoding, canonical form, \texttt{extend\_right},
    \texttt{extend\_left}.  No I/O, no external dependencies.
  \item[\texttt{jellyfish/}] \texttt{KmerDatabase} trait; pure-Rust \texttt{.jf}
    file reader; conditional FFI to libjellyfish (compiled only when
    \texttt{pkg-config} finds jellyfish).
  \item[\texttt{sequence/}] Target FASTA loading with \texttt{needletail}~\cite{needletail};
    \texttt{RefSeq} decomposition into ordered k-mer list; anchor uniqueness validation.
  \item[\texttt{walker/}] Iterative DFS forward and backward walking;
    adaptive threshold computation; bidirectional merge.
  \item[\texttt{graph/}] \texttt{KmerGraph} construction from walk results;
    Dijkstra pathfinding; reference-edge removal; all-shortest-paths enumeration;
    DOT visualisation output.
  \item[\texttt{variant/}] \texttt{VariantType} enum; \texttt{diff\_path\_without\_overlap}
    classifier; NNLS quantifier with bootstrap CI; INDEL normalisation; clustering
    for compound variants.
  \item[\texttt{filter/}] Multi-stage filter pipeline; PoN database; adaptive
    threshold computation; context annotation.
  \item[\texttt{output/}] TSV, CSV, VCF, JSON/JSONL, and XLSX writers.
  \item[\texttt{config/}] TOML configuration loading; \texttt{DetectConfig},
    \texttt{FilterConfig}, etc.
  \item[\texttt{cli/}] Clap derive subcommand definitions for \texttt{detect},
    \texttt{filter}, \texttt{merge}, \texttt{stats}, \texttt{plot},
    \texttt{coverage}, \texttt{run}, \texttt{benchmark}, \texttt{pon}.
  \item[\texttt{confidence/}] Statistical scoring: binomial p-values, Fisher's
    method, Phred-scaled QUAL, strand bias.
  \item[\texttt{trace/}] Detection trace output (per-target graph visualisation,
    walking debug output) for algorithm diagnostics.
\end{description}

\subsection{Type-Driven Design}
\label{sec:type-driven}

A central design principle is that Rust's type system encodes domain knowledge,
making incorrect states unrepresentable.

\textbf{The \texttt{Kmer} newtype} prevents accidentally using a raw \texttt{u64}
as a k-mer.  Operations that are semantically meaningful for k-mers (canonical
form, extension, string conversion) live on the \texttt{Kmer} type; raw integer
arithmetic does not.

\textbf{The \texttt{VariantType} enum} has exactly six variants, matching the
biological classification space.  Code that handles a \texttt{VariantType} must
be exhaustive (the Rust compiler enforces this), preventing missed cases when
new variant types are added.

\textbf{The \texttt{VariantCall} struct} encodes the full variant record with
typed fields: \texttt{rvaf: f64}, \texttt{min\_coverage: u64},
\texttt{pvalue: Option<f64>}.  The \texttt{Option} type makes it impossible
to accidentally access the p-value of a call for which no statistical test was
run; the code \emph{must} handle the \texttt{None} case.

\textbf{The \texttt{WalkResult} struct} wraps both the discovered k-mer node set
and the reference k-mer set together, ensuring they are always passed as a unit
and cannot become inconsistent.

\subsection{Trait Abstractions: KmerDatabase}
\label{sec:trait-abstractions}

The \texttt{KmerDatabase} trait is the central abstraction that allows the same
algorithmic code to run against both the real jellyfish backend and a mock
database for testing:

\begin{lstlisting}[style=rust,
  caption={The KmerDatabase trait (src/jellyfish/mod.rs).}]
pub trait KmerDatabase: Send + Sync {
    fn query(&self, kmer: &str) -> u64;
    fn kmer_length(&self) -> u8;
    // Default: canonicalize then query
    fn query_canonical(&self, kmer: &str) -> u64 {
        self.query(kmer)
    }
}
\end{lstlisting}

The \texttt{Send + Sync} bound is critical: it allows the trait object to be
shared across Rayon threads safely.  \texttt{Send} means the value can be
transferred to another thread; \texttt{Sync} means it can be accessed from
multiple threads simultaneously.  The jellyfish reader satisfies both because
it is read-only after construction.

The \texttt{MockDb} used in tests implements \texttt{KmerDatabase} with an
in-memory \texttt{HashMap<String, u64>}, supporting programmatic insertion of
k-mer counts for precise test control without requiring a real jellyfish
installation.

\subsection{Pure-Rust Jellyfish Reader}
\label{sec:jf-reader}

The original \km{} tool requires the jellyfish Python bindings, introducing a
complex build dependency.  \kmerdet{} provides two backends:

\textbf{FFI backend}: When the build system finds jellyfish via \texttt{pkg-config},
\texttt{build.rs} generates C++ FFI bindings using \texttt{bindgen} and sets the
\texttt{has\_jellyfish} configuration flag.  This backend provides exact
compatibility with all jellyfish versions through the official C API.

\textbf{Pure Rust reader}: The default backend when jellyfish is not installed.
It memory-maps the \texttt{.jf} file with \texttt{memmap2}, parses the JSON
header (k-mer length, hash table size, canonical flag, counter width), and
implements jellyfish's reversible hash function and linear-probing lookup.  The
memory-mapped approach enables $O(1)$ random-access k-mer queries with minimal
RAM overhead.

The conditional compilation is declared in \texttt{build.rs}:

\begin{lstlisting}[style=rust,
  caption={Conditional jellyfish FFI in build.rs.}]
fn main() {
    println!("cargo::rustc-check-cfg=cfg(has_jellyfish)");
    if pkg_config::probe_library("jellyfish-2.0").is_ok() {
        println!("cargo::rustc-cfg=has_jellyfish");
        // compile jf_wrapper/jf_wrapper.cpp
    }
}
\end{lstlisting}

All tests pass without jellyfish installed via the \texttt{MockDb}.

\subsection{Parallelisation Strategy}
\label{sec:parallelism}

Target-level parallelism is the primary source of speedup.  Each target is an
independent unit of work: its walking, graph construction, pathfinding,
classification, and quantification steps depend only on the read-only jellyfish
database.  Rayon's \texttt{par\_iter()} distributes targets across a thread pool
with work-stealing:

\begin{lstlisting}[style=rust,
  caption={Parallel target processing (conceptual, from cli/detect.rs).}]
use rayon::prelude::*;

let calls: Vec<VariantCall> = targets
    .par_iter()
    .flat_map(|target| {
        let ref_kmers = target.decompose(k);
        let walk = walker::walk_bidirectional(&db, &ref_kmers, &cfg);
        let graph = graph::build(&walk);
        let paths = graph::find_paths(&graph);
        paths.iter()
            .filter_map(|path| classify_and_quantify(target, path, &walk, &db))
            .collect::<Vec<_>>()
    })
    .collect();
\end{lstlisting}

The jellyfish database handle is wrapped in \texttt{Arc<dyn KmerDatabase>} so
the reference-counted pointer can be cloned into each thread.  Because
\texttt{KmerDatabase: Sync}, all threads can call \texttt{query()} concurrently
without locks.

The thread pool size defaults to the number of logical CPU cores and is
configurable via \texttt{--threads N}.  On a typical 16-core server, this
provides a $10\text{--}14\times$ speedup over single-threaded execution for a
50-target panel.

\subsection{Error Handling Philosophy}
\label{sec:error-handling}

\kmerdet{} follows a two-level error handling strategy:

\textbf{Library level}: Typed errors using \texttt{thiserror} with variants for
\texttt{InvalidKmer}, \texttt{DatabaseError}, \texttt{WalkingError},
\texttt{ClassificationError}, etc.  These propagate through the library as
\texttt{Result<T, KmerdetError>}, allowing callers to handle specific failure
modes.

\textbf{CLI level}: Errors are converted to \texttt{anyhow::Error} at the
subcommand entry point, which automatically provides contextualised error
messages with the \texttt{.context()}  chain.  Input validation (file existence,
format checks, parameter ranges) occurs eagerly before any computation begins,
so partial output is never written.

\textbf{Per-target errors are non-fatal}: If walking fails for one target
(e.g., non-unique anchor, empty result set), the error is logged at the
\texttt{warn} level and processing continues with the remaining targets.  This
matches \km{}'s behaviour and prevents a single problematic target from blocking
the entire analysis.

\subsection{CLI Design}
\label{sec:cli}

The CLI uses Clap's derive API~\cite{clap}, which generates the argument parser
from annotated structs.  Subcommands map directly to clinical workflow steps:

\begin{lstlisting}[style=bash,
  caption={kmerdet subcommand overview.}]
kmerdet detect   --jf sample.jf --targets panel/ --output calls.tsv
kmerdet filter   --calls calls.tsv --reference known_mutations.tsv
kmerdet merge    --input run1.tsv run2.tsv --output merged.tsv
kmerdet stats    --calls filtered.tsv
kmerdet plot     --calls filtered.tsv --output-dir figures/
kmerdet coverage --jf sample.jf --targets panel/
kmerdet pon      build --controls ctrl1.tsv ctrl2.tsv --output normals.pon
kmerdet run      --jf sample.jf --targets panel/ --reference mutations.tsv
\end{lstlisting}

Default parameters for \texttt{detect} match the \km{} liquid biopsy
configuration: \texttt{--count 2 --ratio 0.05}, replicating \km{}'s
behaviour on identical inputs.  All parameters can be overridden via command
line or a TOML configuration file (\texttt{--config kmerdet.toml}).

Global flags apply to all subcommands: \texttt{--threads N},
\texttt{--verbose/-v} (repeatable), \texttt{--quiet/-q}, \texttt{--no-header}.

\subsection{Testing Strategy}
\label{sec:testing}

\kmerdet{} follows a three-tier testing strategy:

\textbf{Unit tests}: Every public function has at least one unit test.
K-mer encoding/decoding roundtrips, canonical form correctness, extension
logic, and classification rules are verified against known examples.
Property-based tests (using \texttt{proptest}) verify that encoding and
decoding are mutual inverses for all $4^k$ possible k-mers up to $k = 20$.

\textbf{Integration tests}: The test suite in \texttt{tests/} uses \texttt{MockDb}
instances to run the full pipeline on synthetic k-mer distributions with
known ground truth.  Tests cover SNV detection, insertion, deletion, ITD,
homopolymer edge cases, and boundary conditions (zero-count k-mers,
single-k-mer targets).

\textbf{Compatibility tests} (marked \texttt{\#[ignore]}, requiring jellyfish):
Run \kmerdet{} and \km{} on the same \texttt{.jf} databases and target FASTA
files, comparing TSV output field-by-field.  These tests verify that
\kmerdet{}'s reference implementation mode produces identical results to \km{}.

\subsection{Performance Characteristics}
\label{sec:performance}

The k-mer walking algorithm is fundamentally memory-bandwidth bound: each
extension step performs one hash table lookup in the jellyfish database.
Memory-mapped \texttt{.jf} files are optimal for this access pattern because
the OS page cache retains recently accessed pages across queries from different
threads.

For typical targeted sequencing data (50 targets, 200~bp each, $k = 31$):

\begin{itemize}
  \item \textbf{Database size}: $\sim$400~MB per jellyfish database (fits in L3
    cache of modern server CPUs for databases loaded from targeted sequencing).
  \item \textbf{Walking}: $\sim$10{,}000 k-mer queries per target; 50 targets =
    500{,}000 total queries.  At $O(1)$ per query, this completes in $<1$ second
    even single-threaded.
  \item \textbf{NNLS}: The contribution matrix is $31 \times 3$ for a typical
    SNV target; solving NNLS is microseconds.  Bootstrap ($B = 1000$) adds
    $\sim$30~ms per target.
  \item \textbf{Total pipeline}: $\sim$0.5--2 minutes for 50 targets on a
    16-core machine with \texttt{--threads 16}, versus $\sim$2 minutes for the
    Python km component of the \kam{} pipeline on 4 threads.
\end{itemize}
